\id{МРНТИ 27.25.19}{}

\begin{header}
\swa{З.А.~Сарсембаева, Г.И.~Аймичева, А.Е.~Назырова}{ЖЕЛІ АРҚЫЛЫ ЖАЛҒАН КОНТЕНТТІ ТАСЫМАЛДАУДЫ ІЗДЕУ ҮШІН ЖЕЛІЛІК ТРАФИКТІ ТАЛДАУ}

З.А.~Сарсембаева\envelope,
Г.И.~Аймичева,
А.Е.~Назырова
\end{header}

\begin{affil}
Л.Н.Гумилев атындағы Еуразия Ұлттық университеті, Астана, Қазақстан

\corrauthor{Корреспондент-автор: 040203650680@enu.kz}
\end{affil}

Цифрлық технологиялар мен жасанды интеллекттің қарқынды дамуы жалған
ақпарат таратудың жаңа мүмкіндіктерін ашты. Қазіргі уақытта жалған
контентті анықтауға бағытталған зерттеулер көбінесе визуалды талдаумен
шектеледі, ал желі деңгейінде оны анықтаудың тиімді әдістері жеткіліксіз
зерттелген. Сондықтан бұл зерттеу жұмысының негізгі мақсаты желілік
трафикті талдау арқылы жалған контентті, соның ішінде дипфейк мазмұнын
анықтаудың тиімді техникалық әдістерін ұсыну болып табылады. Бұл жұмыста
HTTP-трафикті талдау, файлдарды автоматты түрде шығару, MesoNet
нейрондық желі моделі арқылы дипфейктерді анықтау әдістері қолданылды.
Зерттеу барысында Scapy, PyTorch және Docker технологиялары
пайдаланылды. Нәтижелер көрсеткендей, ұсынылған әдіс дипфейк-контентті
желілік трафик деңгейінде нақты уақыт режимінде жоғары дәлдікпен анықтай
алады. Жүйе қауіпті контентті автоматты түрде анықтап, Telegram арқылы
жедел хабарландыру жібереді, бұл ақпараттық қауіпсіздікті қамтамасыз
етуде үлкен практикалық мәнге ие.

{\bfseries Түйін сөздер:} жалған контент, дипфейктер, ақпараттық
қауіпсіздік, желілік трафикті талдау, нейрондық желілер, Scapy, MesoNet,
автоматты талдау, дезинформация, цифрлық криминалистика.

\begin{header}
АНАЛИЗ СЕТЕВОГО ТРАФИКА ДЛЯ ПОИСКА РАСПРОСТРАНЕНИЯ ЛОЖНОГО КОНТЕНТА ЧЕРЕЗ СЕТЬ

З.А.~Сарсембаева\envelope,
Г.И.~Аймичева,
А.Е.~Назырова
\end{header}

\begin{affil}
Евразийский национальный университет им. Л.Н. Гумилева, Астана, Казахстан,

e-mail: 040203650680@enu.kz
\end{affil}

Стремительное развитие цифровых технологий и искусственного интеллекта
открыло новые возможности для распространения ложной информации.
Особенно в последние годы широко распространяется дипфейк-контент, не
являющийся реальным, но визуально убедительный. В настоящее время
исследования, направленные на выявление ложного контента, в основном
ограничиваются анализом на визуальном уровне, тогда как методы его
эффективного обнаружения на сетевом уровне изучены недостаточно. В связи
с этим основной целью исследования является предложение эффективных
технических подходов к выявлению ложного контента, включая дипфейки,
посредством анализа сетевого трафика. В работе использовались методы
анализа HTTP-трафика, автоматического извлечения файлов и определения
дипфейков с помощью нейронной сети MesoNet. В ходе исследования
применялись технологии Scapy, PyTorch и Docker. Результаты показали, что
предложенный метод способен выявлять дипфейк-контент на уровне сетевого
трафика в режиме реального времени с высокой точностью. Система
автоматически определяет опасный контент и оперативно отправляет
уведомления через Telegram, что имеет большое практическое значение для
обеспечения информационной безопасности.

{\bfseries Ключевые слова:} ложный контент, дипфейки, информационная
безопасность, анализ сетевого трафика, нейронные сети, Scapy, MesoNet,
автоматический анализ, дезинформация, цифровая криминалистика.

\begin{header}
NETWORK TRAFFIC ANALYSIS FOR DETECTING THE DISTRIBUTION OF FALSE CONTENT OVER THE NETWORK

Z.A.~Sarsembayeva\envelope,
G.I.~Aimicheva,
A.Ye.~Nazyrova
\end{header}

\begin{affil}
L.N. Gumilyov Eurasian National University, Astana, Kazakhstan,

e-mail: 040203650680@enu.kz
\end{affil}

The rapid development of digital technologies and artificial
intelligence has opened up new opportunities for the dissemination of
false information. In recent years, deepfake content, which is not
authentic but visually convincing, has become widespread. Current
studies on detecting false content are mostly limited to visual-level
analysis, while effective methods for identifying it at the network
level remain insufficiently explored. Therefore, the main objective of
this research is to propose effective technical approaches for detecting
false content, including deepfakes, through network traffic analysis.
This study applied HTTP traffic analysis, automatic file extraction, and
deepfake detection using the MesoNet neural network model. The research
utilized Scapy, PyTorch, and Docker technologies. The results
demonstrated that the proposed method can accurately detect deepfake
content at the network traffic level in real time. The system
automatically identifies malicious content and promptly sends alerts via
Telegram, which has significant practical value for ensuring information
security.

{\bfseries Keywords:} false content, deepfakes, information security,
network traffic analysis, neural networks, Scapy, MesoNet, automated
analysis, disinformation, digital forensics.

\begin{multicols}{2}
{\bfseries Кіріспе.} Қазіргі таңда жасанды интеллект пен ақпараттық
технологиялардың дамуы арқасында жалған ақпараттың цифрлық кеңістікте
таралу жылдамдығы мен күрделілігі күннен күнге артуда. Жалған контент
тек әлеметтік желілерде ғана емес, сонымен қатар -- форумдар,
хабарламалар жүйесі және заңды сайттар арқылы да таралуда. Бүгінде
жалған ақпарат тек мәтін немесе бейне түрінде ғана емес, мультимедиялық
пішінде, яғни суреттер мен аудиофайлдар арқылы да берілуі мүмкін. Көзге
шынайы көрінетін, бірақ шың мәнінде жалған болып келетін және адам
өміріне, беделіне нұқсан келтіретін контентті алдын ала анықтап, оның
таралуын тоқтату қазіргі кездегі ақпараттық қаіпсіздіктің басты
мәселелерінің бірі.

Жалған контентті танудағы зерттеулер көбіне визуалды ерекшеліктерді
анықтауға бағытталған, ал оларды ерте кезеңде анықтап, оның таралуына
жол бермеу туралы зерттеулер қазіргі кезде аса көп қарастырылмаған және
де тапшылық етуде. Бұл ақпараттық қауіпсіздік саласындағы маңызды
кемшіліктердің бірі.

Осы себептен, бұл зерттеу жұмысының басты мақсаты -- жалған контентті
желілік трафикті талдау арқылы анықтау әдістерін ұсыну. Негізгі зерттеу
сұрағы: желілік трафикті талдау арқылы жалған контенттің таралуын қалай
тоқтатуға болады?

Әдебиетке шолу. Ақпараттық кеңістіктегі жалған контенттің таралуы
қоғамдағы пікірді бұрмалап, халықтың шешім қабылдауына ықпал етеді.
Соңғы зерттеулер көрсеткендей, психологиялық тұрғыдан әсерлі жалған
контент адамдар арасында шыңайы ақпаратқа қарағанда 70\%-ға жылдам
таралады.

Google trends (\url{https://trends.google.com/trends/}) деректеріне
сәйкес, «Fake news» (жалған жаңалықтар), «deepfakes» (жасанды бейнелер),
«disinformation» (әдейі таралған жалған ақпарат) және «misinformation»
(білместіктен таралған жалған ақпарат) сияқты сұраныстар адамдар
арасында тұрақты қызығушылық тудыруда. Мұндай қызығушылықты әсіресі,
саяси, әлеуметтік және экономикалық оқиғалар кезінде байқауға болады.
Бұл халықтың жалған ақпаратқа алданып қалуына деген қорқынышын және
дезинформация қаупын түсінуін көрсетеді.

Соңғы бірнеше жыл ішінде бұл терминдер бойынша іздеу белсенділігі
жаһандық деңгейде ғана емес, Қазақстанда да айқын өсуде. Бұд құбылыс
дезинформацияның тек халықаралық ауқымдағы мәселе екенін көрсетеді. Осы
орайда, жалған контентпен күресуде тиімді техникалық құралдар мен
әдістерді әзірлеу қажеттілігі барған сайын айқындала түсуде.
\end{multicols}

\fig[0.7\textwidth]{i/image2}[1-сурет. Google Trends мәліметтері бойынша жалған ақпарат және
deepfakes тақырыптарына қызығушылық]

\fig[0.7\textwidth]{i/image3}[2-сурет. Дипфейктер мен жасанды интеллект тақырыбы бойынша
ғылыми жарияланымдардың елдер бойынша белсенділігі]

\begin{multicols}{2}
Бұл зерттеу тақырыбы ғылыми ортада да кең танымалдылыққа ие. SciVal
платформасының бибилеометриялық талдау деректеріне сүйінетін болсақ,
2022-2024 жылдар арасында дипфейк, терең оқыту және жасанды интеллект
туралы мақалалар саны бұрынғымен қарағанда әрдеқайда артып келуде.
Әсіресе бұл мәселеге Қытай, Үндістан, АҚШ, Италия және Ұлыбритания
елдерінің ғалымдары ерекше көңіл бөлген. Жарияланымдар саны бойынша
Қытай (644), Үндістан (374) және АҚШ (278) елдері көш бастап тұр. Ал
Италия (85) және Ұлыбритания (83) елдері бір-бірінен қалыспай, алғашқы
бестікті түйіндейді (сурет-2).

{\bfseries Материалдар мен әдістер.} Желілік трафикті талдауға арналған
NTMA әдістемесі арқылы жалған контентті анықтау. Желілік трафикті талдау
-- ақпараттық қауіпсіздікті қамтамасыз ететін басты құралдардың бірі.
Трафикті талдау арқылы әдеттен тыс нәрселерді (анномалияларды), заңсыз
әрекеттерді, фишингтік шабуылдарды және жалған ақпаратты тарату
арналарын анықтауға болады. Соңғы жылдары әлеуметтік желі арқылы жалған
ақпараттын көбеюі және де зиянды бағдарламалардың таралуының өсуіне
байланысты бұл талдаудың ғылыми тұрғыдағы маңызы арта түсуде. Желілік
құрылымдардың күрделенуі және деректерді шифрлау әдістерінің жетілуі
дәстүрлі тексеру құралдарының қолдану аясын шектеп, жаңа тәсілдерді
талап етуде. Сондықтан, машиналық оқыту мен терең оқытуға негізделген
әдістерді қолдану қажетттілігі артуда.

NTMA (Network Traffic Monitoring and Analysis), яғни желілік трафикті
бақылау және талдау -- бұл желінің қозғалысын ең ұсақ бөлшектерге дейін
бақылауға мүмкіндік беретін әдістер жиынтығы. Бұл тәсіл желінің жұмысын
және өнімділігін терең зерттеуге, сондай-ақ пайдаланушылардың желідегі
әрекеттерін талдауға мүмкіндік береді {[}1{]}. Коммуникациялық жүйелер
мен желілер тұрғысынан алғанда NTMA өте маңызды рөл атқарады. Ол
желілердің қалай жұмыс істейтінін түсінуге, олардың өнімділігін
бақылауға, ресурстарды тиімді пайдалануды, оңтайландыруды және
телекоммуникациялық инфрақұрылымды SLA (Service Level Agreement --
қызмет көрсету деңгейі туралы келісім) талаптарына сәйкес басқаруға
көмектеседі.

\fig[0.8\columnwidth]{i/image4}[3-сурет. NTMA процедурасының құрылымы]

Машиналық оқыту мен терең оқыту әдістерін пайдалану арқылы NTMA жүйелері
желілік аномалияларды автоматты түрде анықтай алады. Терең оқыту (Deep
Learning) белгілерді автоматты түрде шығарып, деректердегі жасырын
заңдылықтарды анықтау үшін қолданылады. Бұл технология үлкен көлемдегі
деректерді талдау отырып, модельдердің жұмыс тиімділігін арттырады және
дәстүрлі машиналық оқыту әдістерінен бір саты жоғары тұрады {[}1{]}.

NTMA желі ресурстарын басқаруда және қауіпсіздікті қамтамасыз етуде
басты құрал болып саналады. Терең оқыту әдістерін пайдалану желі
трафигін талдау міндеттерін қазіргі коммуникациялық жүйелердің жоғары
жылдамдық пен сенімділік талаптарына сай шешуге көмектеседі.

Машиналық оқыту шифрланған трафикті талдауға NTMA-ны бейімдеуде маңызды
рөл атқарады. Әдеттен тыс жағдайларды (аномалияларды) анықтау үшін
көбінесе автоэнкодерлер пайдаланылады, олар жүйенің қалыпты
міңез-құлқынан ауытқуын бірден байқайды. Сондай-ақ, K-орташа кластерлеу
сияқты бақылаусыз әдістерге негізделген кластерлеу трафиктің жаңа
үлгілерін табуға ат салысады. Мысалы, автоэнкодерлерді шифрланған
ағындарды тексеруге қолдану HTTPS арқылы деректерді жасырын шығары
әрекеттерін анықтауға мүмкіндік берді {[}1,2{]}.

Желілік трафикті талдауда Deep Packet Inspection (DPI) және NetFlow
әдістерінің үлесі зор. Бұл әдістердің әрқайсысы желідегі ақпарат
ағындарын зерттеуде өзіндік ерекшеліктермен көзге түскен.

Deep Packet Inspection (DPI) - бұл желі арқылы өтетін пакеттердің ішкі
мазмұнын толықтай тексеруге мүмкіндік беретін технология. Яғни,
қарапайым сөзбен айтқанда, ол тек пакеттің тақырыбын ғана емес, сонымен
қатар оның ішкі құрылымында талдайды. DPI-дің басты артықшылықтарының
бірі - желідегі заңсыз әрекеттерді, шифрланған трафиктегі күдікті
үлгілерді және зиянды бағдарламаларды анықтау қабілеті. Бұл әдіс
қауіпсіздік жүйелерінде, атап айтқанда IDS (Intrusion Detection Systems
- шабуылды анықтау жүйелері) және IPS (Intrusion Prevention Systems -
шабуылды алдын алу жүйелері) белсенді қолданылады {[}2{]}.

NetFlow - желілік трафиктің құрылымдық сипаттамаларын талдауға арналған
Cisco компаниясы әзірлеген технология. Ол аномалияларды табу үшін трафик
ағындарының статистикалық мәліметтерін жинайды. NetFlow - DDos
шабуылдарын, ботнеттердің белсенділігін, және фишингтік шабуылдарды
айқындауда тиімді болып табылады.

DPI мен NetFlow тәсілдерін бірге қолдану арқасында желілік қауіпсіздікті
қамтамасыз ету оң нәтижелер береді. DPI арқылы пакеттердің ішкі
мазмұнына назар аударсақ, NetFlow арқылы трафик ағындарының құрылымдық
параметрлерін талдап, желідегі күдікті мінез-құлықты анықтауға болады.
Осы тәсілдердің бірігуі желідегі заңсыз әрекеттерді, зиянды
бағдарламаларды және жалған ақпаратты тарату арналарының белсенділігін
жоғары дәлдікпен анықтауға мүмкіндік береді {[}3{]}.

Көпмодальды тәсілдер дегеніміз -- әртүрлі ақпарат түрлерін (мәтін,
сурет, бейне, аудио) біріктіріп талдауға негізделген әдістер.
Көпмодальды әдістерді пайдалану арқылы жалған ақпаратты таратудың
бірнеше негізгі бағыты пайда болған. Біріншіден, әлеуметтік желілер мен
жаңалықтарда тек мәтіндік контент ғана емес, сонымен қатар суреттер мен
бейнелер белсенді қолданылады. Мұндай медиалық элементтер жалған
ақпараттың шыңайлылық деңгейін күшейту мақсатында жасалады {[}4{]}.

Көпмодальды тәсілдердің желілік криминалистикадағы маңызы зор. Себебі,
жалған ақпарат таратушылар көбіне тек бір ғана дереккөзбен шектелмей,
бірнеше арнаны бір мезгілде қолданады. Мысалы, әлеуметтік желілердегі
жалған жаңалықтар көбінесе мәтін, сурет және бейне форматтарында қатар
беріледі және олардың арасында өзара логикалық байланыс болады. Осы
байланыстарды толыққанды талдау үшін көпмодальды әдістерді қолданамыз.

Жалған ақпаратты көпмодальды әдістер арқылы анықтауда машиналық оқыту
мен терең оқыту әдістері белсенді қолданылады. Бұл әдістер әртүрлі
дереккөздерден алынған ақпаратты біріктіруге және олардың арасындағы
айырмашылықтарды анықтауға мүмкіндік береді {[}5{]}.

Бұл бағыттағы негізгі модельдердің бірі - GNN (Graph Neural Networks -
графтық нейрондық желілер). GNN ақпараттың таралу жолдарын зерттеп,
күмәнді байланыстарды анықтауға көмектеседі. Бұл модельден бөлек,
визуалды контентті өңдеу үшін CNN (Convolutional Neural Netwoorks --
конволюциялық нейрондық желілер) және мәтіндік және аудиофайлдарды
талдау үшін LSTM (Long Short-Term Memory ұзақ мерзімді жады) модельдері
қолданылады. Жалған контентті талдауда қолданылатын әдістердің
тиімділігі мен тәжірибелік маңызын негіздеу мақсатында {[}5-7{]}
әдебиеттердегі ғылыми мақалаларға талдау жасалды, нәтижесі 1-кестеде
ұсынылған.

Көпмодальды тәчілдерді қолдану кезінде ерекше назар аударуға тұрарлық
тағы бір аспект - кросс-модальды сәйкессіздіктерді анықтау. Бұл
операциялар әртүрлі ақпарат көздері арасындағы сәйкестікті бағалауға
мүмкіндік береді. Мысалы, хабарлама мәтіні бір оқиға туралы, бірақ оған
тіркелген бейне басқа оқиға туралы болуы мүмкін. Мұндай айырмашылықтарды
автоматты түрде анықтау үшін көпмодальды нейрондық желіліер қолданылады
{[}8{]}.
\end{multicols}

\tcap{1-кесте. Көпмодальды терең оқыту архитектураларының сипаттамасы}
\begin{longtblr}[
  label = none,
  entry = none,
]{
  width = \linewidth,
  colspec = {Q[106]Q[308]Q[269]Q[256]},
  cells = {font = \footnotesize},
  row{1} = {c,font=\bfseries\footnotesize},
  hlines,
  vlines,
}
Архитектура                           & Сипаттамасы                                                                                                                                                           & Қолдану салалары                                                                                                                                   & Жалған контентті анықтаудағы өзектілігі                                                                                       \\
Көпқабатты персептрон (MLP)           & Тығыз қабаттардан тұрады, әр қабат деректердің жаңа түрін береді.                                                                                                     & Мәтінді, суретті және дыбысты біріктіру үшін қолданылатын негізгі әдіс.Әр түрлі дереккөздерді синхрондау және олардан алынған ақпаратты біріктіру. & Жалған немесе бұрмаланған мазмұнды көрсететін метадеректердегі ауытқуларды анықтау                                            \\
Конволюциялық нейрондық желілер (CNN) & Уақыт қатарлары немесе кеңістіктік деректер сияқты тор деректер құрылымдарын өңдеуге арналған, ол жергілікті үлгілерді (мысалы, кластер өлшемі, ұзақтығы) бөлектейді. & Суреттер мен бейнелерді өндеу, бейнелерді мәтінмен байланыстыру және мазмұн үйлесімділігін тексеру.                                                & Жасанды мазмұнды жеткізуге байланысты орындалатын пакеттік үлгілер мен пайдалы жүктеме шабуылдарында оғаш ауытқуларды анықтау \\
Рекуррентті нейрондық желілер (RNN)   & Деректерді ретімен өңдейді. RNN алдыңғы кірістердің «есту қабілетін» сақтайды, уақыт реттілігін талдау үшін қолайлы етеді.                                            & Мәтіндік және дыбыстық деректерді уақытша байланыстар бойыншы зерттеу, жалған хабарламаларды сұрыптау                                              & Деректер ағынындағы уақытша сәйкессіздіктерді қадағалау арқылы жалған немесе өзгертілген мазмұнды анықтайды                   \\
Ұзақ мерзімді жады (LSTM)             & RNN-ге ұқсас, градиент жоғалуымен байланысты мәселелерді шешеді. Тізбектердегі ұзақ мерзімді маусымдық тәуелділіктерге мамандандырылған.                              & Уақыт бойынша деректердің өзгеруін талдау, мәтін мен аудио контент арасындағы байланысты зерттеу                                                   & Уақыт өте келе өзгертілген немесе жалған ақпаратқа тән ауытқуды анықтау.                                                      \\
Гейттелген рекуррентті блоктар (GRU)  & LSTM-нің жеңілдетілген нұсқасы, есептеу ресурстарын аз талап етеді, бірақ тізбекті деректерді өңдеуде ұқсас қабілеттерге ие                                           & Мәтін, аудио және бейнелерді LSTM-ге қарағанда тезірек талдайды.                                                                                   & Жалған мазмұнды таратудағы трафик үлгілеріндегі қарапайым айырмашылықтарды анықтайды.                                         \\
Автоэнкодерлер                        & Бақылаусыз оқыту үшін қолайлы архитектура. Кіріс деректерін қысылған көріністерге кодтауға және оларды қайта құруға үйретілген нейрондық желілер                      & Әртүрлі дереккөздерден алынған ақпаратты қысу және олардың арасындағы заңдылықтарды тану                                                           & Қайта құрастырылған деректердің бұрмалануымен байланысты сәйкессіздіктер мен ауытқуларды анықтау                              \\
Терең генеративті модельдер           & Вариациялық автоэнкодерлер (VAE) және генеративті қарсылас желілер (GAN) сияқты әдістерді қамтиды.                                                                    & Генеративті модельдерді қолданып, дипфейкті құрады және оны анықтау үшін кері талдау жүргізеді.                                                    & Жалған трафикті имитациялау немесе сенімділікті тексеру үшін қарсы нүктелерді жасау үшін үлгілерді үйрету                     \\
Гибридті модельдер (CNN-RNN)          & CNN кеңістік ерекшеліктерін және RNN уақытша модельдеуін біріктіреді.                                                                                                 & Жалған контенттің сәйкестігін сурет, мәтін және аудионы өңдеу арқылы тексереді                                                                     & Кеңістіктік және уақыттық белгілерді біріктіру арқылы жалған контенттің таралуындағы күрделі заңдылықтарды анықтау            \\
Ансамбльдік модельдер                 & Бірнеше архитектураның артықшылықтарын пайдалану үшін, оларды бір модельге біріктіреді.                                                                               & Мәтін,сурет және дыбыс арасындағы үйлесімділікті айқындау үшін бірнеше үлгілердің нәтижелерін біріктіру                                            & Жалған контентті анықтау үшін бірнеше тәсілдерді біріктіріп қолданады.                                                        
\end{longtblr}

\begin{multicols}{2}
Көпмодальды тәсілдердің тағы бір маңызды артықшылығы - статистикалық
талдау мен мазмұнды бағалау үйлесімі. Бұл тәсіл мазмұнды тарату
динамикасын және оның пайдаланушыға әсерін зерттеуге мүмкіндік береді.
Әлеуметтік желілердегі тенденцияларды талдай отырып, жалған ақпараттың
таралу жылдамдығын және оның аудиториясын анықтауға болады. Бұл талдау
тәсілі жалған ақпараттың табиғатын зерттеуге және оның таралуын алдын
алуға мүмкіндік береді {[}9{]}.

{\bfseries Нәтижелер және талқылау.} Эксперименттік жұмыстың мақсаты -
жалған контентті анықтау үшін желілік трафикті талдауға арналған кешенді
шешім құру. Осы мақсатқа жету үшін келесі міндеттер қойылды:

1. HTTP-трафигін сенімді түрде ұстап қалу;

2. Трафиктен автоматты түрде медиафайлдарды бөліп алу;

3. Таңдалған нейрондық желі арқыла медиаконтентті талдау;

4. Күдікті медиафайл табылған жағдайда, Telegram bot арқылы хабарлама
жіберу.

Бұл күрделі шешімді жүзеге асыру үшін Scapy, PyTorch және Docker
технологиялары таңдалған болатын.
\end{multicols}

\fig[0.5\textwidth]{i/image5}[4-Сурет. Жүйенің жалпы құрылымы]

\begin{multicols}{2}
Жалған контентті анықтау үшін медиафайлдардағы өте ұсақ бөлшектерді
талдай алатын конволюциялық нейрондық желіні қолданатын MesoNet модельі
тандалған болатын. Бұл модельдің басты ерекшелігі - аз есептеу қуатын
қамтамасыз етеді және жалған контентті анықтау дәлдігі жоғары.

\emph{Қадам 1. Ортаны дайындау}. Жүйенің кез келген жерде бірдей және
сенімді жұмыс жасауын қамтамасыз ету үшін Docker технологиясы
қолданылды.

Практикалық жұмыста барлық қажет кітапханаларды тұрақты түрде қолдану
мақсатында dockerfile ресми python:3.8 бейнесіне негізделген болатын.
Dockerfile үздік әзірлеу тәжірибелерін ескере отырып жасалған. Мұнда
pip-ті жаңарту, барлық қажетті кітапханаларды орнату (requirements.txt),
негізгі кодтарды (sniffer.py, deepfake\_detector.py және т.б)
контейнерге көшіру сияқты кезеңдер жазылған.
\end{multicols}

\fig{i/image6}[5-сурет. Dockerfile құрылымы]

\begin{multicols}{2}
\emph{Қадам 2. Қажетті құралдар мен кітапханалар тізімі -
requirements.txt}

Жүйе жұмысын қамтамасыз ету үшін қажетті кітапханалар тізімі жасалған
болатын. Requirements.txt тізімі келесі құралдар мен кітапханаларды
қамтиды:

1. Scapy - желілік трафикті талдауға арналған python кітапханасы.

2. PyTorch - терең машиналық оқыту модельдерін құру және жаттықтыруға
арналған кітапхана (CPU-есептеулеріне ыңғайлы нұсқасы көрсетілген).

3. Requests - Telegram bot API-мен байланысу үшін.

4. OpenCV - суреттер мен бейнелерді өңдеуге арналған кітапхана.
\end{multicols}

\fig{i/image7}[6-сурет. Requirements.txt мазмұны]

\emph{Қадам 3. Практикалық іске асыру: негізгі скриптің жұмыс принципі.}

Практикалық іске асыру python коды арқылы жүзеге асырылады.

\tcap{Листинг 1. HTTP трафигін ұстап қалу және толық жауапты жинау}
{\footnotesize
\begin{verbatim}
from scapy.all import sniff, TCP, Raw, IP
def packet_callback(pkt):
    if not (pkt.haslayer(TCP) and pkt.haslayer(Raw) and pkt.haslayer(IP)):
        return
    ip_layer = pkt[IP]
    tcp_layer = pkt[TCP]
    if tcp_layer.sport != PORT: return
    payload = bytes(pkt[Raw].load)
    conn_key = (ip_layer.src, tcp_layer.sport, ip_layer.dst, tcp_layer.dport)
    if conn_key not in responses:
        if payload.startswith(b"HTTP/1.") and b" 200 " in payload[:15]:
            responses[conn_key] = payload
    else:
        responses[conn_key] += payload
        buf = responses[conn_key]
        if b"\r\n\r\n" in buf:
            head_end = buf.find(b"\r\n\r\n")
            headers_raw = buf[:head_end]
            for line in headers_raw.decode("latin-1").split("\r\n"):
                if line.lower().startswith("content-length:"):
                    content_length = int(line.split(":",1)[1].strip())
                    break
            body_size = len(buf) - (head_end+4)
            if content_length and body_size >= content_length:
                process_full_response(buf, conn_key)
\end{verbatim}
}

Бұл код бөлігі HTTP жауаптарын толық жинауға арналған. Ол әрбір TCP
пакетін өңдеп, жауаптың толық денесін жинайды.

\tcap{Листинг 2. Суретті сақтап, дипфейкке тексеру}
{\footnotesize
\begin{verbatim}
def save_and_analyze_image(data_bytes: bytes, content_type: str):
    ext = ".jpg" if "jpeg" in content_type else ".png"
    filename = f"sniffed_{file_counter}{ext}"
    file_path = os.path.join(SAVE_DIR, filename)
    with open(file_path, "wb") as f:
        f.write(data_bytes)

    is_fake = deepfake_detector.is_deepfake(file_path)
    if is_fake:
        notify_telegram(f"Deepfake detected in file: {filename}")
\end{verbatim}
}

\begin{multicols}{2}
Бұл функция екі маңызды істі атқарады. Біріншіден, суретті жеке файл
ретінде сақтайды. Екіншіден, сақталған суретті deepfake\_detector
модульі арқылы тексереді. Сурет дипфейк деп таңылған жағдайда, Telegram
арқылы хабар жіберіледі.
\end{multicols}

\tcap{Листинг 3. MesoNet модельі}
{\footnotesize
\begin{verbatim}
def is_deepfake(image_path):
    image = Image.open(image_path).convert('RGB')
    image = mesonet_data_transforms['val'](image)
    image = image.unsqueeze(0).to(device)
    with torch.no_grad():
        output = model(image)
    _, pred = torch.max(output.data, 1)
    return pred.item() == 1
\end{verbatim}
}

MesoNet моделі медиафайлды дипфейк деп анықтаған жағдайда true мәнін
қайтарады {[}10{]}.

\tcap{Листинг 4. Telegram арқылы хабарлама жіберу}
{\footnotesize
\begin{verbatim}
def notify_telegram(message: str):
    url = f"https://api.telegram.org/bot{TELEGRAM_TOKEN}/sendMessage"
    requests.post(
        url,
        data={
            "chat_id": TELEGRAM_CHAT_ID,
            "text": message
        },
        timeout=5
    )
\end{verbatim}
}

\emph{Қадам 4. Docker контейнерін іске қосу}

\begin{multicols}{2}
Жасалған снифферді docker ішінде қосу үшін желіге толық қолжетімділік
беру қажет. Сондықтан -\/-network=host және -\/-cap-add=NET\_ADMIN
опцияларын қолданамыз. Бұл Scapy-ге http трафигін ешқандай кедергісіз
өңдеуге мүмкіндік береді.
\end{multicols}

\fig{i/image12}[7-сурет. Контейнерді желіге қосу конфигурациясы]

\emph{Қадам 5. Дипфейк анықталған кезде Telegram арқылы хабарлама
жіберу}

\begin{multicols}{2}
Дипфейк анықталған кезде жүйе автоматты түрде Telegram боты арқылы
хабарлама жіберді. Бұл механизм бізге қауіп-қатерге мүмкіндігінше тез
жауап беруге және шешім қабылдауға мүмкіндік берді.
\end{multicols}

\fig{i/image13}[8-сурет. Телеграм хабарламасы]

\begin{multicols}{2}
Эксперимент нәтижелері біз таңдаған әдістің тиімділігін толық растады:
трафикті ұстап қалудан бастап хабарлама жіберуге дейінгі барлық қадамдар
қатесіз және жылдам орындалды.

{\bfseries Қорытынды.} Бұл зерттеу жұмысы қойылған мақсаттарға жетіп,
жалған ақпаратпен күресте өз үлесін қосуға және ақпараттық қауіпсіздікті
қамтамасыз етуде тиімді қолдана алатынын көрсетті. Scapy, PyTorch,
Docker технологияларын қолдану арқасында желілік трафикті талдап,
жасанды контентті тауып алуға арналған автоматтандырылған кешенді жүйе
құрылды. Жоба әзірге тек http трафигімен ғана шектелгенімен, болашақта
басқа хаттамалармен де интеграцияланады және талдау дәлдігін арттырады
деп үміттенеміз. Жасанды интеллектпен машиналық оқытудың алдағы
жетістіктері бұл шешімнің тиімділігін арттырып, әртүрлі салаларда
қолданылуына мүмкіндік береді. Зерртеу жұмысының барысында құрылған жүйе
болашақта SOC және CERT сияқты құрылымдарда пайдалануға бейімделген.
Алдағы зерттеулер желілік трафикті тереңірек талдауды, шифрланған
протоколдармен жұмыс істеуді қамтамасыз етуі керек. Бұл бағыттағы
зерттеу жұмыстарының артуының арқасында дезинформациямен күресте жаңа
деңгейге шығуға мүмкіндік пайда болады.
\end{multicols}

\begin{center}
{\bfseries Әдебиеттер}
\end{center}

\begin{refs}
1. Abbasi M., Shahraki A., Taherkordi A. Deep learning for network
traffic monitoring and analysis (NTMA): A survey// Computer
Communications. -2021. -Vol.170. -P.19-41. DOI
10.1016/j.comcom.2021.01.021.

2. Bayat N., Jackson W., Liu D. Deep learning for network traffic
classification// arXiv preprint. -2021. -No. arXiv:2106.12693. DOI:
10.48550/arXiv.2106.12693.

3. Alanazi S., Asif S. Exploring deepfake technology: creation,
consequences and countermeasures// Human-Intelligent Systems
Integration. -2024. -Vol.6.-P.49--60. DOI 10.1007/s42454-024-00054-8.

4. Shen Y., et al. Fake news detection on social networks: a survey //
Applied Sciences. -2023. -Vol.13(21): 11877. DOI
10.3390/app132111877.

5. Sattarov O., Choi J. Detection of rumors and their sources in social
networks: a comprehensive survey// IEEE Transactions on Big
Data.-2024.-Vol.11(3) DOI 10.1109/TBDATA.2024.3522801.

6. Gambín Á. F., Yazidi A., Vasilakos A., et al. Deepfakes: current and
future trends// Artificial Intelligence Review. -2024. -Vol.57:64.
DOI 10.1007/s10462-023-10679-x.

7. Meel P., Raj C., Bhawna. A review of web infodemic analysis and
detection trends across multi-modalities using deep neural network //
International Journal of Data Science and Analytics. -2025. -Vol.
20(5). -P.4149--4176. DOI 10.1007/s41060-025-00727-w.

8. Zannettou S., et al. The web of false information: Rumors, fake news,
hoaxes, clickbait, and various other shenanigans// Journal of Data and
Information Quality (JDIQ). -2019. -Vol.11(3). -P.1-37. DOI
10.1145/3309699.

9. Alwhbi I. A., Zou C. C., Alharbi R. N. Encrypted network traffic
analysis and classification utilizing machine learning// Sensors.
-2024. -Vol.24(11): 3509. DOI 10.3390/s24113509.

10. Liu Honggu. MesoNet-Pytorch. -2019. URL:
\url{https://github.com/HongguLiu/MesoNet-Pytorch/tree/master}.
-Қаралған күні: 19.09.2025.
\end{refs}

\begin{center}
{\bfseries References}
\end{center}

\begin{refs}
1. Abbasi M., Shahraki A., Taherkordi A. Deep learning for network
traffic monitoring and analysis (NTMA): A survey// Computer
Communications. -2021. -Vol.170. -P.19-41. DOI
10.1016/j.comcom.2021.01.021.

2. Bayat N., Jackson W., Liu D. Deep learning for network traffic
classification// arXiv preprint. -2021. -No. arXiv:2106.12693. DOI
10.48550/arXiv.2106.12693.

3. Alanazi S., Asif S. Exploring deepfake technology: creation,
consequences and countermeasures// Human-Intelligent Systems
Integration. -2024. -Vol.6.-P.49-60. DOI 10.1007/s42454-024-00054-8.

4. Shen Y., et al. Fake news detection on social networks: a survey //
Applied Sciences. -2023. -Vol.13(21):11877. DOI 10.3390/app132111877.

5. Sattarov O., Choi J. Detection of rumors and their sources in social
networks: a comprehensive survey// IEEE Transactions on Big Data
-2024.-Vol.11(3) DOI 10.1109/TBDATA.2024.3522801.

6. Gambín Á. F., Yazidi A., Vasilakos A., et al. Deepfakes: current and
future trends// Artificial Intelligence Review. -2024. -Vol.57: 64.
DOI 10.1007/s10462-023-10679-x.

7. Meel P., Raj C., Bhawna. A review of web infodemic analysis and
detection trends across multi-modalities using deep neural network //
International Journal of Data Science and Analytics. -2025. -Vol.
20(5). -P.4149-4176. DOI 10.1007/s41060-025-00727-w.

8. Zannettou S., et al. The web of false information: Rumors, fake news,
hoaxes, clickbait, and various other shenanigans// Journal of Data and
Information Quality (JDIQ). -2019. -Vol.11(3). -P.1-37. DOI
10.1145/3309699.

9. Alwhbi I. A., Zou C. C., Alharbi R. N. Encrypted network traffic
analysis and classification utilizing machine learning// Sensors.
-2024. -Vol.24(11): 3509. DOI 10.3390/s24113509.

10. Liu Honggu. MesoNet-Pytorch. -2019. URL:
\url{https://github.com/HongguLiu/MesoNet-Pytorch/tree/master}.
-Қаралған күні: 19.09.2025.
\end{refs}

\begin{info}
\hspace{1em}\emph{{\bfseries Авторлар туралы мәлімет}}

Сарсембаева З. - магистрант, Л.Н. Гумилев атындағы Еуразия ұлттық
университеті, Астана, Қазақстан,
e-mail:~040203650680@enu.kz;

Аймичева Г. - PhD, Л.Н. Гумилев атындағы Еуразия ұлттық университеті,
Астана, Қазақстан, e-mail:
aimicheva\_gi\_1@enu.kz;

Назырова А. - PhD, Л.Н. Гумилев атындағы Еуразия ұлттық
университеті,Астана, Қазақстан, e-mail: nazyrova\_aye\_1@enu.kz.

\hspace{1em}\emph{{\bfseries Information about the authors}}

Sarsembayeva Z.-- Master's student, L.N. Gumilyov Eurasian National
University, Astana, Kazakhstan, e-mail: 040203650680@enu.kz;

Aimicheva G. - PhD, senior, L.N. Gumilyov Eurasian National University,
Astana, Kazakhstan, е-mail:
aimicheva\_gi\_1@enu.kz;

Nazyrova A. - PhD, senior lecturer , L.N. Gumilyov Eurasian National
University, Astana, Kazakhstan, е-mail: nazyrova\_aye\_1@enu.kz.
\end{info}
